\chapter{文献综述与现有商业方案的信息边界}
\label{ch:lit}

沿着\ref{ch:intro}章给出的那条链，本章处理第一环之前的准备工作：已有的文献与货架方案各自闭合了哪一段，还留下哪一段。智能蒸烤一体机同时落在食品工程、逆热传导和安全控制三条文献线上。食品工程给出体积平均、等温线和品质动力学；逆热传导给出由表面恢复内部热阻的病态结构；安全控制给出屏障投影和家电型式试验的测量点。接到四个执行器上之后，真正进入积分器的只有若干层：体积平均给出能量、液相和蒸汽方程的骨架，品质动力学把停机规定为进度阈值，逆热传导说明稳态观测对内部热阻灵敏度为零，对偶控制允许同一阀门在某一窗口内为辨识服务，屏障函数把参考动作投影到由硬测量定义的安全集。工程核已经运行的是十六条控制体、标量跟踪器和 \(150\)--\(190\,\mathrm{s}\) 的固定时间窗。交叉表比较的，就是这四条策略在同一积分器上的停机偏差。

商品目录把这三条文献线压缩成一页彩页。容积、上下管功率、蒸汽发生量、风机档位和“智能菜单数量”写在参数表里。参数表闭合的是边界能力：加热管能把腔温送到多少度，蒸汽阀一分钟能送多少克水，风机能把换热系数抬到哪一档。食品内部的热阻、冷点位置和这一次过程里的进度，不在参数表的任何一行。穿刺针作为选配件出现时，说明书写的是“插入食品中心”，没有写冷点可以偏离几何中心，也没有写针尖改变局部场。图像识别作为卖点出现时，说明书写的是“自动匹配菜单”，蒸汽遮挡使画面中断的那几秒不在文案里。文献综述因此不只是把论文排队，而是把商品路径已经闭合的量和仍未进入后验的量分开。

\section{食品热湿传递的描述框架}

Whitaker 给出干燥过程中多孔介质里热量、质量和动量同时传递的体积平均理论\cite{whitaker1977}。固相骨架、液相和气相分别平均，界面交换记为源项。Datta 把这一骨架用于食品加工中的同时热质传递\cite{datta2007}：蒸发和冷凝是界面源，蒸汽作为相变项进入能量方程。积分器沿用平均之后能量、液相和蒸汽方程的结构，把蒸对应为冷凝潜热为正的源，把干壳对应为近表面渗透率与有效导热下降。干燥器或某一烘箱的边界闭合、文献中的物料数组，都留在各自的测定条件里；本次核用自己的界面通量和常数。

Feyissa、Gernaey 与 Adler-Nissen 给出对流烤箱中肉块的三维耦合场，并表明冷点可以偏离几何中心\cite{feyissa2013}。冷点因此定义为场的本质下确界，在一维核里取十六条控制体的最低温度。一块 \(24\) 乘 \(16\) 的显式截面把脂肪列放在下标 \(8\)，\(20\,\mathrm{min}\) 后最低温度仍在绝热侧，偏离中心 \(13.75\,\mathrm{mm}\)。该截面复现“冷点可以不在中心”，没有复现三维烤箱里的冷点云。

Choi 与 Okos 给出组分热物性的加和传统\cite{choi1986}。有效导热可以随含水变化，各向异性不在该加和之内。本次核用常数 \(\alpha=2.0\times 10^{-7}\,\mathrm{m}^2/\mathrm{s}\) 和 \(\rho c=3.4\times 10^{6}\,\mathrm{J}/(\mathrm{m}^3\mathrm{K})\) 给出基线 \(k=0.68\,\mathrm{W}/(\mathrm{m\cdot K})\)，脂肪夹层把指定控制体乘 \(0.45\)。

水分活度采用 GAB 三参数形式\cite{gab1980}
\begin{equation}
\label{eq:gab}
\frac{M}{X_m}=\frac{C K a_w}{(1-K a_w)(1-K a_w+C K a_w)}.
\end{equation}
公开拟合表\cite{sorption2026}中诊断码为 \(0\) 的清洁子集有 \(140\) 条，类别以谷物、果蔬和粉体为主，肉类条数为 \(0\)。土豆清洁拟合的温度只到 \(60^\circ\mathrm{C}\)。该表约束淀粉类活度在测定温度内的数量级。烤制表面温度高于测定上限，肉糜也没有对应行。核里的活度是含水的单调饱和函数，与 \eqref{eq:gab} 同方向、不同反演：含水升高则活度升高并趋近 \(1\)。图\ref{fig:gab}画出土豆六条、番茄温度族以及面粉、大米、小米、奶粉、蘑菇对照。正文只列土豆短表与十二条代表性记录，完整 \(140\) 条见附录。

\begin{figure}[htbp]
\centering
\colimg{S_gab.png}
\caption{公开 GAB 清洁子集的等温线。左：土豆吸附与解吸；右：番茄温度族与谷物、粉体对照}
\label{fig:gab}
\end{figure}

清洁子集上 \(X_m\) 的十分位、中位、九十分位为 \(0.0366\)、\(0.0733\)、\(0.1453\)；\(C\) 为 \(3.98\)、\(11.20\)、\(45.24\)；\(K\) 为 \(0.607\)、\(0.773\)、\(0.959\)。分位按 140 条诊断码为 0 的 GAB2 记录线性插值。表中单层容量 \(X_m\) 为 \(0.03\sim 0.22\) 量级的小数，能量常数 \(C\) 大于 1，列顺序与源表一致。

\begin{table}[htbp]
\centering
\caption{土豆清洁拟合，测定温度只到 \(60^\circ\mathrm{C}\)}
\label{tab:potato}
\setlength{\tabcolsep}{3.2pt}
{\small
\begin{tabular}{llrrrrrr}
\toprule
记录 & 方向 & \(T/^\circ\mathrm{C}\) & \(X_m\) & \(C\) & \(K\) & \(n\) & \(R^2\) \\
\midrule
c074 & ADS & 30 & 0.0567 & 14.34 & 0.912 & 7 & 0.9963 \\
c075 & ADS & 45 & 0.0457 & 13.71 & 0.885 & 7 & 0.9970 \\
c076 & ADS & 60 & 0.0359 & 6.38 & 0.762 & 7 & 0.9990 \\
c077 & DES & 30 & 0.0763 & 9.21 & 0.861 & 7 & 0.9932 \\
c078 & DES & 45 & 0.0549 & 8.55 & 0.860 & 7 & 0.9975 \\
c079 & DES & 60 & 0.0522 & 3.75 & 0.666 & 7 & 0.9984 \\
\bottomrule
\end{tabular}}
\end{table}

\begin{table}[htbp]
\centering
\caption{GAB2 清洁子集的代表性记录}
\label{tab:gab12}
\setlength{\tabcolsep}{3.0pt}
{\small
\begin{tabular}{llrrrrr}
\toprule
记录 & 类别 & \(T\) & \(X_m\) & \(C\) & \(K\) & \(R^2\) \\
\midrule
c012 & milkPowder & 5 & 0.0720 & 4.21 & 0.859 & 0.988 \\
c017 & milkPowder & 25 & 0.0375 & 7.19 & 0.944 & 0.999 \\
c028 & tomato & 20 & 0.2037 & 27.48 & 0.834 & 0.998 \\
c031 & tomato & 30 & 0.1621 & 22.65 & 0.897 & 0.985 \\
c041 & tomato & 60 & 0.1165 & 9.36 & 0.960 & 0.983 \\
c047 & beetroot & 50 & 0.1169 & 1.65 & 0.926 & 0.990 \\
c074 & potato & 30 & 0.0567 & 14.34 & 0.912 & 0.996 \\
c077 & potato & 30 & 0.0763 & 9.21 & 0.861 & 0.993 \\
c081 & rice & 5 & 0.0697 & 5.85 & 0.787 & 0.995 \\
c088 & mushroom & 30 & 0.1722 & 4.72 & 0.713 & 0.999 \\
c098 & flour & 10 & 0.0823 & 15.13 & 0.722 & 0.999 \\
c101 & millet & 20 & 0.0743 & 12.82 & 0.796 & 0.994 \\
\bottomrule
\end{tabular}}
\end{table}

清洁子集以谷物、果蔬和粉体为主，肉类名称的条数为 0。这些行约束淀粉类活度在测定温度内的数量级。完整 140 条见附录。


\section{热力进度作为停机依据}

van Boekel 把食品品质表述为温度指数、转化率幂和含水开关的乘积\cite{vanboekel2008}。品质是反应进度：同一条中心温度曲线可以对应不同的表面褐变和不同的热力进度，停机因此规定为进度阈值，温度阈值只作为跟踪器的中间量。

Holdsworth 把热力安全进度表述为沿冷点温度历史的等效时间积分\cite{holdsworth1997}。评分核沿用
\[
\frac{\mathrm{d}F}{\mathrm{d}t}=10^{(T_{\mathrm{cold}}-T_{\mathrm{ref}})/z},\qquad
\alpha_m=1-\exp(-F/F_{\mathrm{req}}),
\]
参考温度 \(70^\circ\mathrm{C}\)，\(z=10\,\mathrm{K}\)，\(F_{\mathrm{req}}=2\,\mathrm{min}\)。设计线 \(\alpha_m=0.85\) 对应 \(F=3.794\,\mathrm{min}\)，用作策略比较的零点。微生物初始数量留在规格之外，交叉表给出的是同一条进度轴上四条策略的相对位置。

家用说明书很少写出 \(F\) 值，更常见的是“中心 \(70^\circ\mathrm{C}\)”“表面金黄”“锁住水分”三句并排，而三句对应三个不同的量：中心温度是跟踪器的中间量，表面金黄是褐变代理 \(\alpha_b\)，锁住水分是表面含水。停机读的是 \(\alpha_m\)，由冷点温度历史积分而来。恒温冷点若保持 \(70^\circ\mathrm{C}\)，走到设计线还需 \(3.794\,\mathrm{min}\)；若只到 \(60^\circ\mathrm{C}\)，同一条线要将近 \(38\,\mathrm{min}\)。厨房里“看起来已经熟了”往往读的是表面褐变和蒸汽量，冷点进度可以仍停在 \(0.1\) 以下。

\section{内部热阻估计的病态性质}

由表面温度反演内部热流是病态的，高频噪声随厚度被放大\cite{ozisik2000}。半无限体上角频率 \(\omega\) 的谐波在深度 \(L\) 处振幅衰减为 \(\mathrm{e}^{-kL}\)，\(k=\sqrt{\omega/(2\alpha)}\)。要从表面恢复该深度的同频分量，测量噪声乘上 \(\mathrm{e}^{kL}\)。\(40\,\mathrm{s}\) 探针窗在 \(30\,\mathrm{mm}\) 处的放大约为 \(1.46\times 10^{8}\)；开环菜单时间量级上放大降到几十，仍不足以把厚件冷点上的谐波读出来。只加密表面红外像素，提高的仍是已经衰减掉的高频分量，冷点要靠穿透时间尺度上的动态激励，或一个对内部热阻敏感的积分量。

集合卡尔曼滤波提供带协方差的状态更新\cite{evensen2003}。完整估计器可以表述为分层集合滤波。本次日志里的估计器是一个标量跟踪器：冷点由表面经时间常数跟上，时间常数正比于热阻估计和厚度的平方。跟踪器与集合滤波分开报告。物理信息网络和神经算子可以加速参数化演化\cite{raissi2019,li2021fno,lu2021deeponet}，外推范围由训练轨迹里的执行器覆盖决定。两节点教师上训练的主干加分支网络留出均方根约 \(41^\circ\mathrm{C}\)，评分循环仍使用标量跟踪器。

\section{对偶控制与安全屏障}

Feldbaum 指出控制输入同时推动状态和辨识\cite{feldbaum1965}。家电执行器只有有限探针库。信息度量采用先验熵与后验熵之差\cite{lindley1956}。结构因果模型里的干预是替换结构方程的输入后再积分；条件相关保留控制策略本身\cite{pearl2009}。加大蒸汽可能只是控制器对当前低温的反应，效应要由干预分布给出。

模型预测控制给出有限时域的执行量\cite{rawlings2017}。控制屏障把参考动作投影到由测量定义的安全集\cite{ames2017cbf}。家用电器对漏电、温升和异常工作的通用要求见 IEC 60335-1 与 GB 4706.1\cite{iec60335,gb4706}。水位、门、剩余电流和加热器温度因此进入安全投影，进度估计留在品质回路；正文给出的是测量隔离的结构。

\section{现有商品方案的信息边界}

商品蒸烤箱已经能把空气温度稳住，已经能按重量调出一条曲线，已经能在门上装一根穿刺针。稳住的是边界，调出的是申报类别，刺中的是针尖所在的那一点。表\ref{tab:stop}区分已经能够稳定调节的边界量和仍然隐藏的食品量。蒸烤一体合并的是执行器：四个边界进入同一腔体，内部状态仍要由这一次过程观察。

\begin{table*}[t]
\centering
\caption{已有路径的信息止点}
\label{tab:stop}
\begin{tabularx}{\textwidth}{l>{\raggedright\arraybackslash}p{4.2cm}>{\raggedright\arraybackslash}p{4.6cm}}
\toprule
路径 & 已闭合的量 & 未进入后验的量 \\
\midrule
机械定时 & 无 & 全部内部状态 \\
腔温闭环 & 空气温度一点 & 冷点、含水、进度 \\
按重量调用曲线 & 申报质量与类别 & 初温、厚度、脂肪、形状 \\
穿刺探针 & 刺入点温度 & 未刺中的冷点；刺入改变局部场 \\
表面红外 & 辐射亮度 & 发射率、反射、内部梯度 \\
腔体湿度闭环 & 相对湿度设定误差 & 食品内部水分与冷凝位置 \\
图像分类后调用曲线 & 外观类别 & 中心进度；蒸汽遮挡使图像中断 \\
\bottomrule
\end{tabularx}
\end{table*}

机械定时的停止时刻是开环函数。厨房里对应的是面板上的分钟旋钮：拧到 \(20\)，铃声响就停。本次核的开环策略按厚度给出菜单时间 \(t=900+18000(L-0.015)\) 秒：\(15\,\mathrm{mm}\) 为 \(900\,\mathrm{s}\)，\(30\,\mathrm{mm}\) 为 \(1170\,\mathrm{s}\)，\(50\,\mathrm{mm}\) 为 \(1530\,\mathrm{s}\)。这是“曲线按厚度档标定、内部热阻从未被这次过程观察”的工程形式。薄件的导热时间尺度在菜单时刻已经走完，菜单把它送到进度饱和；厚件的导热时间尺度尚未走完，菜单在冷点起来之前就切断加热。

腔温闭环闭合空气温度，不估计冷点，也不停机。面板上对应的是温度旋钮：拧到 \(180^\circ\mathrm{C}\)，管壁和热风把腔温拉住，铃声仍由分钟旋钮决定，或者干脆一直加热到有人开门。识别食材后调用固定曲线，曲线不随这一次的后验更新。图像分类把外观对应为类别，蒸汽遮挡使图像中断，中心进度仍不在画面里。穿刺能降低停机不确定度，但改变局部场，而且不刺中冷点时读数不代表 \(T_{\mathrm{cold}}\)。一根针从门缝伸进去，刺中的是针尖所在的那一毫米，脂肪夹层和几何偏心可以把真正的冷点推到另一侧。穿刺因此只作为对照通道，主观测向量不依赖它。

上述路径闭合的都是边界量、申报类别或刺入点，没有把这一次过程里的内部热阻写入后验。用腔温或湿度误差调节执行器，停机时刻仍然不看见冷点。把传感器曲线映射成三维着色，也不改变停止判据。商品一体机把蒸、烤、组合写在面板上，内部仍是四条允许集合：蒸档夹住加热上限，烤档关掉蒸汽并抬高排汽，组合档放开四个量。允许集合改变通量的上界，冷点进度仍由厚度和这一次过程里看见的热阻决定。

把货架方案按“已经装上的传感器”与“闭合的量”排成一张表，比按品牌排更能看出止点在哪里。表\ref{tab:product}列出五类典型方案。它们共同的形状是：传感器装在腔体侧或刺入点，闭合的是边界温度、边界湿度或一个点的读数；没有一类把“这一次过程里的内部热阻”写进后验，因为那需要的不是更贵的传感器，而是一段离开平衡的激励。

\begin{table*}[t]
\centering
\caption{货架方案按传感器与闭合量排列。共同止点是“内部热阻不进入后验”}
\label{tab:product}
\begin{tabularx}{\textwidth}{l>{\raggedright\arraybackslash}p{3.4cm}>{\raggedright\arraybackslash}p{3.4cm}>{\raggedright\arraybackslash}p{3.6cm}}
\toprule
方案类型 & 已装传感器 & 闭合的量 & 与本文的差别 \\
\midrule
机械定时 $+$ 腔温闭环 & 腔温热电偶、定时器 & 空气温度与时间 & 时间写在厚度档上，冷点从未被观测 \\
按重量调用曲线 & 称重托盘 & 申报质量与类别 & 初温、厚度、脂肪夹层不改写曲线 \\
穿刺探针 & 针尖温度或湿度 & 刺入点读数 & 未刺中冷点；针尖改变局部场 \\
表面红外 $+$ 图像分类 & 红外阵列、摄像头 & 表面亮度与外观类别 & 发射率未知；蒸汽遮挡使画面中断 \\
边界湿度闭环 & 湿度探头或导出管 & 腔体相对湿度 & 食品内部水分与冷凝位置仍隐藏 \\
\midrule
本文 & 表面温度 $+$ 平均含水 $+$ 腔温 & 冷点进度后验（含热阻） & 150--190 s 主动探针购买第二个灵敏度方向 \\
\bottomrule
\end{tabularx}
\end{table*}

\section{公开专利的信息止点}

课程要求把文献、专利与产品一并对照。表\ref{tab:patent}列出四项国内主流家电企业的公开方案。美的探针装置把可伸缩针携带温度或湿度传感器伸入食品\cite{cn110793657}，闭合的仍是刺入点读数，未刺中的冷点与针尖对局部场的扰动都不进入后验。美的湿度检测方法经冷却管导出腔体气体后再测相对湿度\cite{cn111110016}，闭合的是边界相对湿度，食品内部水分与冷凝位置仍隐藏。方太蒸箱烤箱一体机采用双腔冷凝蒸发结构\cite{cn109965678}，闭合蒸汽排放与冷凝水回收，执行器接口仍是加热、蒸汽、排汽与风机四个量。老板电器的蒸汽冷凝组件用相变材料承接排气侧蒸汽\cite{cn114886312}，闭合的是排气侧冷凝，内部热阻同样不进入后验。

\begin{table*}[t]
\centering
\caption{国内主流家电企业公开专利的信息止点}
\label{tab:patent}
\begin{tabularx}{\textwidth}{l>{\raggedright\arraybackslash}p{3.4cm}>{\raggedright\arraybackslash}p{3.4cm}>{\raggedright\arraybackslash}p{3.6cm}}
\toprule
专利 & 权利人 & 已闭合的量 & 未进入后验的量 \\
\midrule
CN110793657B & 美的厨房电器 & 刺入点温度或湿度 & 未刺中的冷点；针尖扰动 \\
CN111110016B & 美的厨房电器 & 导出气体的相对湿度 & 食品内部水分与冷凝位置 \\
CN109965678B & 方太厨具 & 双腔冷凝蒸发与一体结构 & 内部热阻与冷点进度 \\
CN114886312A & 老板电器 & 排气侧相变冷凝 & 内部热阻与表面结壳 \\
\bottomrule
\end{tabularx}
\end{table*}

四项专利覆盖探针测温、腔体湿度、蒸烤一体结构和蒸汽冷凝，代表目前货架方案的技术止点：边界闭环或单点穿刺。本报告的技术路线落在非平衡主动探针与硬测量安全隔离上：短时窗给内部热阻第二个灵敏度方向，硬测量把品质估计从断电条件里拿出去。

\section{菜单时钟对内部状态的替代}

打开一台家用蒸烤一体机的彩页，第一眼落在容积、上下管功率、蒸汽发生量和菜单数量上。容积告诉操作者能放几层烤盘，功率告诉管壁能把腔温送到哪一档，发生量告诉阀侧一分钟能送出多少克水，菜单数量告诉面板上能翻出多少条名字。四行数字闭合的都是边界能力。一块从冰箱取出的肉糜进门之后，冷点进度仍按半厚平方走。彩页很少另开一栏写内部热阻。穿刺针作为选配件出现时，插图把针尖画在几何中心，而二维截面把最低温度落在绝热侧，偏移 \(13.75\,\mathrm{mm}\)；摄像头作为卖点出现时，画面被蒸汽窗里的白雾盖住，中心进度仍不在镜头里。

菜单表把食材名字对应到停机时刻。欧式面包、清蒸鲈鱼、烤肉，三条名字对应三条时钟；重量感应托盘若装上了，申报质量只用来换一行，初温、脂肪夹层和真实厚度仍不改写那一行的分钟数。货架把名字写得越细，厨房越容易把名字当成内部模型；积分器把名字压缩为活度尺度，停止判据仍读冷点进度。

一体机相对分体布置省下的是台面。蒸箱和烤箱并进同一块钢板之后，冷凝和干热必须共用四个旋钮。为冷点多送一股蒸汽，薄件上是冷凝收益，厚件上可能只把表面打湿；为表面排湿而加大排汽，冷点又可能失去冷凝潜热。符号冲突落在界面通量上。彩页把冲突压缩成一句“蒸烤随心切换”；厨房里切换的是允许集合，冷点能不能起来仍由 Fourier 数决定。

\section{安全切断与模型外附件}

彩页很少单独画加水口。家用一体机的水箱多半藏在侧舱或底部抽屉，容量约 \(1\,\mathrm{L}\)，加水口旁边常有一扇窄刻度窗。注满之后浮子抬起，干簧管允许加热和蒸汽离开零；水箱空了仍加热是干烧。厨房里更常见的一次停机来自水箱没有推到位：浮子没有抬起，面板给出“请加水”，四个量全部被固定在零。积分器把这一条对应为水位测量，不进入品质跟踪器。

搁架常见五档。上层靠近加热管，辐射份额更大，表面先上色；下层靠近蒸汽喷嘴的冷凝区，表面更容易被打湿。核把搁架高度压缩进 \(h=12+55u_{\mathrm{fan}}\) 的同一条对流，没有单独的层高状态。烤盘冲孔改底部换热，冷凝托盘接住门玻璃内侧流下来的水。操作者抽托盘时门会开几毫米，微动落下，加热和蒸汽被固定在零。

面板上的“预热完成再放菜”对应的是食品初温。核从食品初温与腔体室温附近一起积分。从冰箱直接进门的肉糜停在低档，回温后面团停在高档。预热把食品从冷藏拿到室温，相当于在格子里换一列初温，菜单分钟数并不跟着改。开环仍按厚度档切断，对偶仍按点估计、方差和厚度停机。预热缩短的是导热时间尺度上的一段，厚件的 Fourier 数在家用时域内仍远小于 \(1\)。

童锁、预约和腔体照明写在说明书的另几页。童锁锁住旋钮，预约把开始时刻往后挪，照明只改可见性，三件事都不进入积分器。除垢会改变阀侧发生量的长期均值，核把发生量对应为 \(u_{\mathrm{st}}\) 的瞬时开度，钙化不单独成状态。

五条文献各自已经成立，接到四个执行器上之后仍缺模块之间的判据。体积平均给出能量、液相和蒸汽方程的骨架，不规定何时停机。品质动力学把停机规定为进度阈值，不规定观测从哪里来。逆热传导说明稳态观测对内部热阻灵敏度为零，但不给出家电上那 \(40\,\mathrm{s}\) 该怎么拧阀门。对偶控制允许同一阀门为辨识服务，却不隔离安全测量。屏障函数把参考动作投影到硬测量定义的安全集，却不估计冷点。缺的是何时后验已经够用，何时必须改变边界以提升内部参数的可辨识性，信息写入哪几个参数，效应沿哪条路径改变品质，以及品质模型失效时谁接管执行器。已经运行的核把这些判据压缩为可以比较的同一套旋钮。
